\documentclass[aspectratio=169,professionalfonts,11pt]{beamer}
\usepackage{luatexja}
\renewcommand{\kanjifamilydefault}{\gtdefault}
\usepackage{amsmath,amssymb,booktabs,tikz}
\usetikzlibrary{arrows.meta,positioning,calc}
\definecolor{lessonaccent}{HTML}{246B95}
\setbeamercolor{structure}{fg=lessonaccent}
\setbeamercolor{normal text}{fg=black,bg=white}
\setbeamertemplate{navigation symbols}{}
\setbeamertemplate{footline}[frame number]
\setbeamersize{text margin left=9mm,text margin right=9mm}
\title{論文を背景から理解する}
\author{}
\date{}
\begin{document}
% Replace the scaffold with source-grounded content. No frame/page limit.
% lesson-frame: title
\begin{frame}\titlepage\end{frame}
% lesson-frame: learning-map
\begin{frame}{学習の道筋}\tableofcontents\end{frame}
\section{背景と前提知識}
% lesson-frame: foundations-start
% 問い・前提・図/具体例・結論・次との接続・出典を記してから制作する。
\begin{frame}{この論文で理解したいこと}
承認した読者設定と、原論文に基づく説明をここから作成します。
\end{frame}
\end{document}
